\section{Completion of the genericity argument}
\label{sec:completion}

We first collect the quantifiers in the finite experiment. This is
necessary because a smooth neighborhood controls a prescribed finite
number of derivatives, whereas the packet construction can demand
arbitrarily high finite accuracy.

\subsection{Finite choices in one experiment}\label{completion:destruction}

Fix a passing test $\mathcal A(O,B)$, a datum
$d\in\mathcal A(O,B)$, and a finite data order $m\ge10$.
We explain why the choices in Proposition~\ref{holo:violation} can
be made in an order that achieves Equation~\eqref{test:density-target}.

By Corollary~\ref{test:branch-seed}, choose $w\in O$ whose first
future runout from the comparison development is a branch. Prepare
the labels so that $v\to\infty$ and $u\to u_*<\infty$, using the
independently prepared exchanged gauge for the other branch. Fix the
strict observation-domain and collar margins about its late points.

The estimates of Sections~\ref{energy:section}--\ref{holo:section}
provide the three strict comparisons consumed by the experiment:
\begin{equation}\label{completion:margins}
\begin{aligned}
 \|H_T^\pm\|_{C^{30}}&=o(\sqrt{a_p}),\\
 e^{o(T)}\left(\|F_T^{\mathrm{odd}}\|_{C^2}
       +\sum_\pm\|H_T^\pm\|_{C^2}^2\right)&=o(A_T),\\
 e^{o(T)}(a_p/h_T+h_T^{16})&=o(A_T),
 \qquad h_T=e^{-\kappa T/16}.
\end{aligned}
\end{equation}
Here $a_p\asymp e^{-\kappa T}$ and $A_T=e^{-7\kappa T/8}$.
The first comparison permits the adaptive observer to return to a
nearby initial seed and controls the connection in its frame.
The second preserves the signed curvature signal,
and the third contradicts that signal if the selected datum still
passes the transport test. All subexponential factors refer to fixed
finite orders.

Fix first the data order $m$ and all the finite output and buffer orders
used in Propositions~\ref{prep:signed-entry}--\ref{holo:violation}.
These include the interior energy, ordinary supremum, odd-source,
constraint, exterior, trace and normal-gauge-jet orders that supply the
metric and loop estimates in Equation~\eqref{completion:margins}.
Choose a total exponential loss smaller than $\kappa/32$ and allocate
it among the finitely many estimates of
Propositions~\ref{prep:signed-entry} and~\ref{energy:signed-experiment}.
As specified there, next fix the bulk thresholds, small auxiliary mass,
late-support fraction, exterior radius and protected-ball tolerance.
Their order makes the small exterior growth and the required early
agreement hold simultaneously.

Now choose $\zeta>0$ sufficiently small that every positive frequency
power at those fixed output orders fits its allocated loss. Only after
this choice increase the finite transport, Taylor and matching degrees
to obtain the required inverse accuracy. For example,
$A_T\lambda_T^{-J}e^{o(T)}\le e^{-4\kappa T}$ follows for large $T$
from a fixed choice with $7\kappa/8+J\zeta>4\kappa$.
Proposition~\ref{bg:packet} ensures that increasing this accuracy does
not increase the positive frequency powers at the chosen output
orders. The additional coefficient orders and constants remain finite;
their subexponential bounds are used after the formal degrees are fixed.
Finally take $T$ sufficiently large for all these estimates, the
prescribed $p_m$ proximity, and membership $w_T^\pm\in O$.

Proposition~\ref{prep:signed-entry} supplies the two exact complete
constraint data with their entire corrected exterior tails, and
Proposition~\ref{energy:signed-experiment} supplies their finite vacuum
evolutions and signed estimates. Proposition~\ref{attach:development}
places the observation neighborhoods in the full developments
$\mathcal M_{d_T^\pm}$; Proposition~\ref{attach:observer} gives the
nearby observing seeds and the uniform comparison with the test norms.
For at least one sign, Proposition~\ref{holo:violation} then produces
a finite loop violating the same $\mathcal A(O,B)$ test. Its base lies
at a positive proper time before the geodesic's full-development
endpoint, and the loop lies in the future of the complete initial
surface. Since $p_m(d_T^\pm-d)\to0$ for both signs, this proves the
remaining density assertion.

\begin{proof}[Proof of Theorem~\ref{phase:main}]
Work first at the fixed normalized center of
Section~\ref{phase:scaling}.
Lemma~\ref{test:closed} makes each passing set
\(\mathcal A(O,B)\) relatively closed in \(\mathcal U\).

It has empty relative interior. Indeed, given
\(d\in\mathcal A(O,B)\) and any relative neighborhood \(V\) of \(d\)
in \(\mathcal U\), Lemma~\ref{phase:baire} gives a sufficiently
small finite \(p_m\) ball contained in \(V\).
The strict \(p_{10}\) inequality defining \(\mathcal U\) has positive
slack at \(d\), so the radius can be chosen to preserve it.
Proposition~\ref{holo:violation} then gives a datum in
that ball outside \(\mathcal A(O,B)\).

There are countably many basic seed sets \(O\) and integers \(B\).
Consequently
\[
 \mathcal G=
 \bigcap_{\substack{O\text{ in the chosen basis}\\ B\in\N,\ B\ge1}}
       \bigl(\mathcal U\setminus\mathcal A(O,B)\bigr)
\]
is a dense \(G_\delta\) in the nonempty Baire space \(\mathcal U\).
By Proposition~\ref{test:exit-cover}, every datum admitting a
future \(C^1\) extension in Definition~\ref{phase:extension}
belongs to at least one passing set. No datum of \(\mathcal G\)
therefore admits such an extension. In particular the conclusion
concerns the full maximal development and every future exit in
the definition, without any vacuum assumption on the ambient
extension. The rescaling in Section~\ref{phase:scaling} proves
the result for arbitrary fixed \(M>0\).
\end{proof}
